\documentclass[11pt]{article}
\usepackage[margin=1in]{geometry}
\usepackage{amsmath,amssymb,mathtools}
\usepackage{fontspec,unicode-math}
\setmainfont{Latin Modern Roman}
\setsansfont{Latin Modern Sans}
\setmathfont{Latin Modern Math}
\usepackage{xurl,hyperref}
\hypersetup{colorlinks=true,urlcolor=blue,linkcolor=blue}
\newcommand{\sref}[2]{\hyperlink{source:#1:#2}{[S#2]}}
\newcommand{\cataloguescope}[1]{\paragraph{Recorded mathematical question.}#1}
\begin{document}
\section{Local metric homogeneity forcing periodic triangulations}\label{statement}

\subsection{Definitions and mathematical statement}

\paragraph{Shared notation.}$\mathbb N=\{1,2,\ldots\}$, and $\mathbb Z,\mathbb Q,\mathbb R,\mathbb C$ denote the integers, rationals, reals and complex numbers. Any other conventions are specified locally.


Let $X$ be either the Euclidean plane with its Euclidean metric or the hyperbolic plane with its standard constant-curvature metric. Let $\mathcal T$ be a locally finite geometric triangulation of $X$ whose triangles have diameter at most a fixed $r>0$. Assume that for every pair of vertices $v,w$ there is an ambient isometry of $X$ mapping $B_X(v,r)$ onto $B_X(w,r)$ and carrying the restriction of $\mathcal T$ in the first ball onto its restriction in the second. Periodicity means that a discrete group of ambient isometries preserving the entire triangulation acts cocompactly on $X$ (in the Euclidean case this gives a lattice of translation periods).
\paragraph{Statement recorded in the supplied source.}
Must $\mathcal T$ be periodic? \sref{statement}{2}
\subsection{Short English statement}
Does identical geometric structure in all vertex-centered balls of the stated radius force a globally periodic triangulation?
\subsection{Sources}
\begin{description}
\item[\hypertarget{source:statement:1}{S1}] ULAM.AI and the UnsolvedMath Contributors, \emph{UnsolvedMath: A Curated Collection of Open Mathematics Problems}, catalogue entry. CC BY 4.0. \url{https://huggingface.co/datasets/ulamai/UnsolvedMath}.
\item[\hypertarget{source:statement:2}{S2}] Itai Benjamini, \emph{Euclidean vs. Graph Metric}, version dated 16 July 2012; original author manuscript preserved in Arquivo.pt capture of 31 December 2020, Section 5, Question 5.7 (with Romain Tessera), p. 7. \url{https://arquivo.pt/noFrame/replay/20201231041538id_/http://www.wisdom.weizmann.ac.il/~itai/erd100.pdf}.
\end{description}
\label{end:statement}
\end{document}
